\section{Analysis and Discussion}

\subsection{Why DPO Fits This Problem}

The corrupted-text setting gives unusually clean preference pairs.
For many failures, the model's wrong answer is directly copied from the corrupted auxiliary text.
Thus the comparison is explicit: the image-grounded answer should be preferred over the text-following wrong answer.
DPO optimizes exactly this comparison.
This complements, rather than invalidates, the SFT mitigation in Words or Vision~\cite{deng2025words}.
SFT with text augmentation teaches the model from a balanced mixture of cases and can reduce broad text reliance.
Our DPO variant instead uses a narrower but sharper signal: for each mined failure, it tells the model which specific text-derived answer should lose to the visual answer.
In contrast, GRPO must first sample useful candidate answers and then rely on a hand-designed reward.
If the sampled candidates rarely contain the correct answer, or if most candidates are similar text-following errors, the relative reward signal is weak.
This explains why GRPO logs can show diagnostic reward events while greedy held-out predictions barely change.

Our DPO construction uses the benchmark answer as the chosen response.
This is appropriate because the task is not open-ended preference alignment but targeted VQA correction: the correct answer should beat an incorrect answer copied from misleading text.
The advantage over SFT is contrastive.
SFT says only what to output, while DPO also identifies the tempting text-derived answer that should be suppressed.

\subsection{Why Not Simply Ignore the Text?}

Ignoring auxiliary text is an appealing but incomplete baseline.
It would avoid some corrupted-text failures, but it does not solve the actual decision problem faced by a VLM.
At inference time, the model is not given a label saying whether the text is matched, irrelevant, or corrupted.
If the model learns to discard text wholesale, it can lose useful information in match cases and may fail to exploit reliable OCR or retrieved context.

Our objective is therefore not ``use less text'' but ``use text conditionally''.
The model should keep using textual context when it is visually supported and reject it when it conflicts with the image.
DPO is appropriate because it changes the model's preference in exactly the conflict cases we can identify:
\begin{equation}
\text{image-grounded answer} \succ \text{text-following wrong answer}.
\end{equation}
This is different from input ablation or prompt-level instruction to ignore text.
The latter can reduce exposure to misleading text, but it does not teach the model which answer should be preferred when both visual evidence and misleading text are present.
Our use of match and irrelevant preservation metrics is intended to check precisely this point: a successful method should improve corrupted robustness without collapsing into text avoidance.

\subsection{What Counts as Solving Text Bias?}

We therefore define a successful correction operationally through four criteria:
held-out corrupted accuracy should improve, incorrect aux-hit should decrease, match/irrelevant settings should be preserved, and the trend should be at least partially reproducible across nearby settings such as DPO hyperparameters or related VQA benchmarks.
Our DocVQA results satisfy the first three criteria strongly and the fourth criterion partially: Qwen-family transfer is positive, while LLaVA and TextVQA reveal clear boundary conditions.

\subsection{Why Multi-Condition Error Mining Helps}

Text-following span DPO is the core filtering strategy, but the best result comes from multi-condition error-mined DPO.
This does not mean that unfiltered errors are sufficient.
Rather, corrupted errors remain strictly filtered for text-following behavior, while other conditions contribute lower-weight stability signals.
This helps prevent the model from over-specializing to the corrupted prompt format.
The gain is modest because non-corrupted failures are sparse in DocVQA, but it consistently improves the main held-out metric.

\subsection{Boundary Conditions}

Our transfer experiments reveal a clear boundary condition.
The method works when the base VLM already has sufficient visual and OCR competence.
Qwen-family models satisfy this condition on DocVQA and improve substantially.
LLaVA-Next models, under our prompt/template, have very low base DocVQA accuracy.
In this regime, DPO cannot create document-reading ability from scratch; it can only alter preferences among answers the model can already represent.

TextVQA experiments lead to a related conclusion.
DocVQA-trained DPO transfers weakly to TextVQA and does not reliably fix LLaVA-specific TextVQA conflict failures.
Even after mining larger LLaVA-specific TextVQA DPO data and adding strict+preserve pairs, corrupted performance does not improve.
This suggests that TextVQA conflict involves different OCR and answer-format dynamics than DocVQA.
Therefore, we treat TextVQA and LLaVA results as boundary analyses rather than main success claims.

\subsection{Limitations}

Our study has several limitations.
First, the text-following filter is operational rather than causal; accidental overlap with corrupted text is possible.
Second, most experiments use one held-out split and limited seeds due to compute constraints.
Third, our GRPO implementation uses free-form sampled answers, and candidate-constrained variants may be stronger.
Finally, the method assumes ground-truth answers for preference construction, making it suitable for supervised mitigation but not directly for fully unlabeled deployment.
